\documentclass[10pt]{article}

\usepackage[utf8]{inputenc}

\usepackage[T1]{fontenc}

\usepackage{amsmath}

\usepackage{amsfonts}

\usepackage{amssymb}

\usepackage[version=4]{mhchem}

\usepackage{stmaryrd}

\usepackage{bbold}

\usepackage{graphicx}

\usepackage[export]{adjustbox}

\graphicspath{ {./images/} }



\begin{document}

Если $A(\cdot)$ - почти периодич. отображение (см. Лuнейная система дифференциальных уравнений $с$ почти периодическими коәффициентами), то для У. х. п. системы (1) необходимо и достаточно, чтобы система (1) была почти приводима.



Для того чтобы характсристич. показатели системы (1) были устойчивы, достаточно, чтобы нашлось JIяпунова преобразовапие, приводящее систему (1) к клеточно-диагональному виду



\[

\left.\begin{array}{c}

\dot{y}_{i}=B_{i}(l) y_{i}, \\

y_{i} \in \mathbb{R}^{k_{i}}, i=1, \ldots, m, \sum_{i=1}^{m} k_{i}=n,

\end{array}\right\}

\]



такому, что: а) клетки интегрально разделешы, т. е. найдутся числа $d>0, a>0$ такие, что



\[

\left\|Y_{i}(\tau, \theta)\right\|^{-1} \geqslant d(\exp [a \cdot(\theta-\tau)])\left\|Y_{i+1}(\theta, \tau)\right\|

\]



для всех $\theta \geqslant \tau \geqslant 0, i=1, \ldots, m-1$ (здесь $Y_{i}(\theta, \tau)-$ Коши оператор систөмы $(2))$; б) верхний и нижниї центральные показатели системы (2) равны друг другу:



$\Omega\left(B_{i}\right)=\omega\left(B_{i}\right)$ при всятом $i=1, \ldots, m$.



Условия этой теоремы являются и необходимыми условиями У. х. п. системы (1) (см. [6]). Системы с неустойчивыми характеристич. показателями могут обладать свойством стохастической У. х. ІІ.



Характеристич. показатели системы (1) наз. с т о$\mathbf{x}$ а т ически уст о й ч и вы ми (или уст ойч и выми по чти на в ерное), если при $\sigma \rightarrow 0$ характеристич. показатели Ляпунова системы



\[

\dot{y}=A(t) y+\sigma^{2} C(t, \omega) y

\]



стремятся с вероятностью 1 к характеристич. показателям Ляпунова системы (1); здесь элементы матрицы, задающей линейные операторы $C(t, \omega): \mathbb{R} n \longrightarrow \mathbb{R}^{n}$ (в нек-ром - но зависящем от $(t, \omega)$-базисе пространства $\mathbb{R}^{n}$ ) суть независимые ненулевые белые шумы.



Если отображкение $A(\cdot): \mathbb{R} \longrightarrow \operatorname{Hom}\left(\mathbb{R}^{n}, \mathbb{R}^{n}\right)$ равномерно неп рерывно и



\[

\sup _{t \in \mathbb{R}}\|A(t)\|<+\infty,

\]



то для почти всякого отображения $\tilde{A}()$, где



\[

\tilde{A}(t)=\lim _{k \rightarrow \infty} A\left(t_{k}+t\right),

\]



характерістич. показатели системы $\dot{x}=\tilde{A}(t) x$ стохастически устойчивы (для сәвигов динамической системь $(S=\operatorname{Hom}(\mathbb{R}, n, \mathbb{R} n)$ рассматривают нормированную инвариантную меру, сосредоточенную на замыкании траектории точки $A(\cdot)$; под почти всяким $\tilde{A}(\cdot)$ понимается почти всякое $\tilde{A}(\cdot)$ в смысле всякой такой меры).



Пусть динамич, система на гладком замкнутом многообразии $V^{n}$ задана гладким векторным полем. Тогда для почти всякой (в смысле всякой нормированной инвариантной меры) точки $x \in V^{n}$ характеристич. показатели системы уравнений в вариациях вдоль траектории точки $x$ стохастически устоӥчивы.



\begin{center}

\includegraphics[max width=\textwidth]{2023_07_09_d72a9f0ed5fd3007f539g-1(1)}

\end{center}



\includegraphics[max width=\textwidth, center]{2023_07_09_d72a9f0ed5fd3007f539g-1}

70; [4] Н е м ы ц к и й В. В., С т е п а но в В. В., Качественная теория дифференциальных уравнений, 2 изд., $1949,[5]$ Б нова и ее приложения к вопросам устойчивости, М., 1966; [6] нова и ее приложения к вопросам устойчивости, М., $1966 ; 6]$

Итоги науки и техники. Математический анализ, т. 12, М., 1974, Птоги науки и техники. Математический анализ, т. 12, М., 1974,

c. $71-146$.

B. М. Милионциков.



УСТРАНИМАЯ ОСОБАЯ ТОЧКА однозначной аналитическої функции $f(z)$ комплексного переменного $z$ - термин Для обозначения такої точки $a$, в шроколотой окрестнос'и $V(a)=\{z \in \mathbb{C}: 0<|z-a|<\delta\} \quad$ к-рой функция $f(z)$ аналитическая и ограниченная. При этих условиях сугествует конечный предел $\lim _{z \rightarrow a, z \neq a} f(z)=f(a)$.



Этот предел иринимают за значение $f(z)$ в точке $a$ и получалот аналитич. функцию во всей окрестности точки $a$. $\quad E$. Д. Соломенцев.



УСТРАНИМОЕ МНОЖЕСТВО $E$ точек гомплексной плоскости $\mathbb{C}$ лля нек-рого класса $K$ однозначных аналитических функций относительно области $G \subset \mathbb{C}-$ такое компактное множество $E \subset G$, что любая функция $f(z)$ класса $K$ в $G \backslash E$ продолжается как функция класса $K$ на всю область $G$. Иначе эту ситуацию описывают словами «множество $E$ с т и р а е т с я для класса $K$ » или «E есть нуль-множество для класса $K »$, сокращенно: $E \in N(K, G)$. Предполагается, что дополнение $G \backslash E$ есть область и что класс $K$ опрөдөлен для любой области.



Согласно другому огределению, множество $E$ у с тр а ним о для класса $K, E \in N(K)$, если из того, тто $f(z)$ есть функция класса $K$ в дополнении $\overline{\mathbb{C}} \backslash E$, следует, что $f(z)=$ const. При этом включения $E \in N(K, G)$ и $E \in N(K)$, вообще говоря, не равносильны.



Первым результатом об У. м. можно считать классич. те о р м у Ћ о IIинимо й ос о б е н ност и: если функция $f(z)$ аналитическая и ограниченная в проколотой окрестности $V(a)=\{z: 0<|z-a|<\delta\}$ точки $a \in \mathbb{C}$, то она продолжается аналитически в точку $a$. Более широкая постановка вопроса (п р о б л е м а П е н л е в е) принадлежит П. Пенлеве (Р. Painlevé): найти условия на множество $E$, необходимые и достаточные для того, чтобы $E \in N(A B, G)$, где $K=A B$ - класс ограниченных аналитич. функций (см. [1]). Сам П. Пенлеве нашел достаточное условие: $E$ должно иметь нулевую линейную меру Хаусдорфа. Необходимое и достаточнос условие в проблеме Пенлеве получено Л. Альфорсом (см. [2]): $E \in N(A B, G)$ тогда и только тогда, когда $E$ имеет нулевую аналитическую емкость. Существует пример множества $E$ положитөльной длины, но нулевой аналитич. емкости. Об У. м. для различных классов аналитич. функций одного комплексного переменного и относящихся к ним нерешенных проблемах cल. [3], [4], [6], [9].



В случае аналитич. функциї $f(z)$ многих комплексных переменных $z=\left(z_{1}, \ldots . z_{n}\right), n \geqslant 2$, постановка задач об У. м. изменяется в силу классич. т ө о р е мы О с г у д а - Б р а ун а: өсли $f(z)$ - регулярная аналитическая в области $G \subset \mathbb{C}^{n}$ функция, за исключением, быть может, компактного множества $E \subset G$ такого, что дополнение $G \backslash E$ связно, то $f(z)$ аналитически продолжается на всю область $G$. Другие теоремы об У. м. при $n \geqslant 2$, а также их связь с понятием голоморфности области см., напр., в [7], [10].



Задачи об У. м. ставятся также для гармонич., субгармонич. функций и др. Так, напр., пусть $G-$ область евклидова пространства $\mathbb{R}^{n}, n \geqslant 3, E-$ компакт, $E \subset G, H B$ - класс ограниченных гармонич. функций, $H D$ - класс гармонит. функций с конечным интегралом Дирихле. Тогда включения $E \in N(H B, G)$ и $E \in N(H D, G)$ равносильны и имеют место тогда и только тогда, когда емкость $E$ равна нулю (см. [5], [8]).



Jum.: [1] Z $z_{0} \mathrm{r}$ e t t $i$ L., Leçons sur le prolongement analytique, P., 1911; [2] A h l f o r's L., "Duke Math. J.», 1947, v. 14, № 1, р. $1-11$, [3] н о с и р о К., Предельные множества, пер. с англ., М., 1963; [4] Х а в и н с о н С. Я., в сб.. Итоги науки. с о н Л., Избранные проблемы теории исключительных множеств, пер. с англ., M., 1971; [6] М е л ь н и к о в М. С., С ин а н я н С. О., в сб.: 'итоги науюи и техники. Современные проблемы математики, т. 4, М., 1975, с. 143-250; [7] III а 6 а Т Х е й м а н У. К., Ке н Н е д и П. Б., Субгармонические функции, пер. с англ., Lт. 1], М., 1980; [9] Д о л ж е н к o E. П., "t a u s J., "Math. Z.», 1978, Bd 158, S. 45-54. E. Д. Соломенцев.





\end{document}